% !TEX root = chapter-standalone.tex

\section{Interfaces}\label{sec:interfaces}

\todotext{This section needs to be rewritten and expanded; in particular more examples need to be added}

One way to get to know  \SY{semicategories} is to understand them as a generalization of semigroups. 

Here is a little story to illustrate how one might think about this. Consider the ropes (with lengths) of~\cref{chap:putting-things-together}, which had this composition rule:
%
\begin{equation}
    \prfperiod{
        \rope{ropecola}{ropecolb}{2}{0}{0}{1}{$a$}
    }{
        \rope{ropecola}{ropecolb}{2}{0}{0}{1}{$b$}
    }{
        \rope{ropecola}{ropecolb}{4}{0}{0}{1}{$a+b$}
    }
\end{equation}
%
If we are working with all possible ropes whose lengths may be any natural number, then we could say that this gives a kind of model of the semigroup $\tupp{\natnumbers,+}$, because setting ropes end-to-end means that the resulting length is the sum of the component rope lengths. In this setting, \emph{any} two ropes may be composed by being set end-to-end. 

Next, consider electrical chords instead of ropes. Each electrical chord has a socket end and a plug end (and, as with the ropes, also a length):
\begin{equation}
\cordG{a}
\end{equation}

Furthermore, worldwide there is a handful of different types of electrical plugs (and sockets). So let's suppose the socket and plug end of our electrical chords are each labeled with a specific type, and that the socket end and the plug end of a given chord don't necessarily have to have the same type (this can be achieved with adapters, for example). For fun, we'll use little country flags as indicators of the types of the plugs and sockets: 
\begin{equation}
    \label{eq:want}
    {\cord{\TypeIrish}{a}{\TypeSwiss}}
\end{equation}

As with the ropes, we can also think about composing electric chords together to make longer ones. However, now we can only connect two chords if the type of the plug of one of them has the type of the socket of the other:
\begin{equation}
    \label{eq:irish-swiss-result}
    \prfperiod{
        \cord{\TypeIrish}{a}{\TypeIrish}
    }{
        \cord{\TypeIrish}{b}{\TypeSwiss}
    }{
        \cord{\TypeIrish}{a + b}{\TypeSwiss}
    }
\end{equation}

Or, phrased in a slightly more abstract manner, for any plug/socket types $\Obja, \Objb, \Objc$ the general formula to compose chords is
%
\begin{equation}
    \label{eq:cords-generic}
    \prfperiod{
        \cord{\Obja}{a}{\Objb}
    }{
        \cord{\Objb}{b}{\Objc}
    }{
        \cord{\Obja}{a+b}{\Objc}
    }
\end{equation}
 
If we think of the ends of the chords as ``interfaces'', then we can say that two chords are only then composable at a pair of interfaces when those interfaces match. Here there are in fact two conditions at play: for interfaces to match we need 1) that one of the interfaces is a plug and the other is a socket, and 2) we need that both interfaces have the same type. 

Both of these conditions are indicative of how things work generally in semicategories: as we will see shortly in more detail, in semicategories we have \emph{morphisms} (the analogue of electrical chords) which each have a \emph{source} and a \emph{target} (analogous to have a socket end and a plug end). Furthermore, every source and target of a morphism has a \emph{type}, and any two morphisms are composable at a pair of interfaces if 1) the one interface is a target and the other is a source and 2) if the types of the interfaces are the same. 

In summary, we first looked at a specific illustrative model of a semigroup given by ropes with length, and then we generalized this picture by adding ``interfaces'' in order to give a first intuition about semicategories. In the very special case when all interfaces happen to have the same type, we are back to having just a semigroup. It is in this sense that semicategories are a generalization of semigroups. 


\todotext{\bernina: Would be nice to continue the example by considering voltages (220 v 110), voltages converters, appliances that can use one or both.}

%
%
%\begin{figure}
%  \begin{center}
%    \begin{tabular}{c@{\hskip 1cm}c@{\hskip 1cm}c@{\hskip 1cm}c@{\hskip 1cm}}
%      \figplug{A} & \figplug{B} & \figplug{C} & \figplug{D} \\
%      \figplug{E} & \figplug{F} & \figplug{G} & \figplug{H} \\
%      \figplug{I} & \figplug{L} & \figplug{M} & \figplug{N}
%    \end{tabular}
%  \end{center}
%
%\end{figure}

\begin{marginfigure}
    \centering
    \figplug{A}

    \figplug{B}

    \figplug{C}

    \figplug{D}

    \figplug{E}

    \figplug{F}

    \figplug{G}

    \figplug{H}

    \figplug{I}

    \figplug{L}

    \figplug{M}

    \figplug{N}
    \caption{Plug/socket types used in the world}
    \label{fig:plugs}
\end{marginfigure}

% \begin{equation}
%   \label{eq:irish-swiss}
%   \prftree{\cord{\TypeIrish}{c}{\TypeIrish}}{\cord{\TypeIrish}{0}{\TypeSwiss}}{\cord{\TypeIrish}{c}{\TypeSwiss}}
% \end{equation}
